% TOP_PROBLEM: {"id": 30006948, "title": "Smooth Realization of Finite-Entropy Ergodic Transformations", "category_id": 11, "category": {"id": 11, "name": "dynamical_systems", "display_name": "Dynamical Systems", "description": "Problems about long-term behavior of deterministic systems, Hamiltonian dynamics, and periodic orbits.", "slug": "dynamical-systems", "order_index": 11, "created_at": "2026-07-31T15:26:25.671Z"}, "status": "open"}
% FIRST_POSTED: 2026-09-27
% !TeX program = lualatex
% ATTEMPT_STATUS: unresolved
\documentclass[11pt]{article}
\usepackage[margin=25mm]{geometry}
\usepackage{fontspec}
\setmainfont{Latin Modern Roman}
\usepackage{amsmath,amssymb,mathtools}
\usepackage{unicode-math}
\setmathfont{Latin Modern Math}
\usepackage{xurl,hyperref}
\hypersetup{colorlinks=true,urlcolor=blue}
\setcounter{secnumdepth}{0}
\setlength{\parindent}{0pt}
\setlength{\parskip}{5pt}
\setlength{\emergencystretch}{3em}
\begin{document}
\section*{\texorpdfstring{Smooth Realization of Finite-Entropy Ergodic Transformations}{Smooth Realization of Finite-Entropy Ergodic Transformations}}\label{30006948}
\subsection{Definitions and mathematical statement}
Let $T$ be an invertible ergodic measure-preserving transformation of a standard nonatomic probability space $(X,\mu)$. Ergodicity means that invariant measurable sets have measure zero or one. Its Kolmogorov--Sinai entropy is the supremum, over finite measurable partitions $\mathcal P$, of
\[
 \lim_{n\to\infty}\frac1n
 H_\mu\!\left(\bigvee_{j=0}^{n-1}T^{-j}\mathcal P\right),
\]
where $H_\mu(\{P_i\})=-\sum_i\mu(P_i)\log\mu(P_i)$.

Assume this entropy is finite. Must there exist a compact finite-dimensional smooth manifold $M$, a smooth probability volume form $\omega$, a $C^\infty$ diffeomorphism $f$ preserving $\omega$, and a measure-space isomorphism $\phi$ modulo null sets with $\phi\circ T=f\circ\phi$ almost everywhere? The manifold and its dimension may depend on $T$. Realization using a singular invariant measure instead of smooth volume would not meet the selected target.
\par\smallskip\textit{Formulation and concept references:} \href{https://arxiv.org/html/2203.10655}{[S1]}.

\subsection{Short English statement}
Can every finite-entropy ergodic invertible system be realized by a smooth volume-preserving diffeomorphism on some compact manifold?
\subsection{Research attempt: entropy bounds, exact models, and the smoothing gap}
\textbf{Status.} A universal smooth realization is not constructed. We prove an elementary entropy obstruction in fixed dimension, classify the connected one-dimensional volume-preserving case, give a positive-entropy smooth model, and exhibit an exact symbolic realization whose failure of smoothness cannot be repaired simply by changing its values on a null set.

\textbf{Source audit.} Open Problem 8 in Appendix B of [S1] asks for a compact manifold and smooth invariant volume, rather than an arbitrary invariant probability on a compact space. Its discussion of presentations also records finite entropy as a necessary condition. The present argument verifies that necessity by a metric covering estimate and does not interpret symbolic realizations or realizations with singular measures as answers to the selected question. The inspected survey and a search for subsequent work did not supply a general realization theorem to invoke here; this is a source-access statement, not a claim to have checked every announcement.

\subsubsection{1. Why finite entropy is necessary}
Let $f$ be a $C^1$ diffeomorphism of a compact $d$-dimensional smooth manifold. Choose a Riemannian metric. There is a finite Lipschitz constant $L\geq1$ for $f$. On a disconnected manifold, there are finitely many components and their mutual distances can be chosen positive; this changes constants but not the argument.

For $n\geq1$, define the orbit metric
\[
 d_n(x,y)=\max_{0\leq j<n}d(f^jx,f^jy).
\]
If $d(x,y)<\varepsilon L^{-(n-1)}$, then $d_n(x,y)<\varepsilon$. Compact smooth $d$-manifolds admit covers by at most $C\delta^{-d}$ balls of radius $\delta$, for sufficiently small $\delta$: this follows by taking finitely many bounded coordinate charts with uniformly comparable Euclidean and Riemannian distances. Therefore an $n$-orbit spanning set at precision $\varepsilon$ has cardinality at most
\[
 C_\varepsilon L^{d(n-1)}.
\]
The definition of topological entropy by spanning sets gives
\[
 h_{\mathrm{top}}(f)\leq d\log L<\infty.                 \tag{1}
\]
Using the standard entropy inequality $h_\omega(f)\leq h_{\mathrm{top}}(f)$ for an invariant Borel probability, every proposed smooth model has finite metric entropy.

Metric entropy is preserved by a measure isomorphism: pull a finite measurable partition back under the isomorphism. Every finite join has exactly the same atom probabilities, so its Shannon entropy is unchanged. Passing to the limit and then taking the supremum over partitions gives
\[
 h_\mu(T)=h_\omega(f).                                  \tag{2}
\]
Consequently infinite entropy is excluded by (1), even if the measure-space isomorphism is completely discontinuous.

On the other hand, (1) does not supply a universal finite upper bound for the problem. Both $d$ and the derivative size of the realizing map are allowed to depend on $T$. One cannot disprove the question by selecting a sufficiently large but finite entropy unless an independent restriction on all possible dimensions and derivatives is proved.

\subsubsection{2. The connected one-dimensional case is rigid}
Consider first a circle with a smooth positive probability density $\rho(x)\,dx$. The cumulative-volume coordinate
\[
 h(x)=\int_0^x\rho(t)\,dt\pmod1
\]
is a smooth circle diffeomorphism and takes this volume to Lebesgue probability. Thus an orientation-preserving volume-preserving diffeomorphism becomes an orientation-preserving Lebesgue-preserving circle diffeomorphism $g$.

For a lift $G:\mathbb R\to\mathbb R$, preservation of the length of every sufficiently short interval implies
\[
 G(b)-G(a)=b-a.
\]
It follows that $G(x)=x+\alpha$, so $g$ is a rigid rotation. In the smooth setting the same conclusion follows from $G'=1$. If preservation is interpreted for the measure while allowing orientation reversal, the lift has derivative $-1$ and the map is $x\mapsto c-x$, an involution. Such an involution is not ergodic on a nonatomic probability circle: a sufficiently small arc and its image have invariant union of measure strictly between zero and one. If preservation is interpreted as $f^*\omega=\omega$ for an oriented volume form, reversal is excluded already by that condition.

The rotation $R_\alpha$ is ergodic exactly when $\alpha$ is irrational modulo one. For irrational $\alpha$, an invariant $L^2$ function with Fourier coefficients $\widehat F(k)$ satisfies
\[
 (e^{2\pi i k\alpha}-1)\widehat F(k)=0,
\]
so all nonconstant coefficients vanish. For rational $\alpha$, unions of short arcs over a finite orbit give nontrivial invariant sets. Every rotation is an isometry, hence has topological and metric entropy zero by the covering argument with $L=1$.

An interval with smooth positive density is likewise put into cumulative-volume coordinates. A measure-preserving increasing diffeomorphism is the identity and a decreasing one is reflection, neither ergodic for nonatomic volume. If a compact one-dimensional manifold has several components, ergodicity forces $f$ to cycle through all components, and an iterate restricts to one of the preceding types. The identity $h(f^k)=k h(f)$ then again gives zero entropy. This identity follows from grouping the joins defining partition entropy into blocks of length $k$ and taking suprema. Thus positive entropy requires dimension at least two, but the problem permits that increase.

This supplies a complete affirmative subcase: any system already measurably isomorphic to an irrational rotation has a smooth volume-preserving realization on the circle, with the original irrational rotation itself as model. No Diophantine restriction appears, since the required conjugacy from the abstract system is only measurable.

\subsubsection{3. A positive-entropy smooth example with elementary checks}
On $\mathbb T^2$ take
\[
 f_A(x)=Ax\pmod{\mathbb Z^2},\qquad
 A=\begin{pmatrix}2&1\\1&1\end{pmatrix}.
\]
The determinant is one and the inverse matrix has integer entries, so $f_A$ is a smooth invertible map preserving Lebesgue area. Its eigenvalues are
\[
 \lambda=\frac{3+\sqrt5}{2}>1,\qquad\lambda^{-1}.
\]
For an invariant $F\in L^2(\mathbb T^2)$, Fourier coefficients are constant along the orbit of each frequency under $A^T$. Every nonzero integer frequency has an infinite orbit: if $(A^T)^qk=k$ for $q>0$, neither eigenvalue of $(A^T)^q$ is one, so $k=0$. Square summability of the coefficients forces their common value on any infinite orbit to be zero. Thus $f_A$ is ergodic.

We can verify positive metric entropy without assuming that this map realizes any particular Bernoulli system. Choose a finite measurable partition $\mathcal P$ whose atoms have diameter less than $\varepsilon$, with
\[
 \varepsilon(1+\|A\|)<1,
\]
and $\varepsilon$ smaller than the injectivity radius of the flat torus. If $x,y$ lie in one atom of $\bigvee_{j=0}^{n-1}f_A^{-j}\mathcal P$, lift each small displacement between $f_A^jx$ and $f_A^jy$ to a vector $\delta_j\in\mathbb R^2$ of norm at most $\varepsilon$. Then
\[
 \delta_{j+1}-A\delta_j\in\mathbb Z^2,
 \qquad |\delta_{j+1}-A\delta_j|<1,
\]
so that integer vector is zero. Therefore $\delta_j=A^j\delta_0$.

Since $A$ is symmetric, its stable and unstable eigenvectors can be chosen orthonormal. The stable coordinate of $\delta_0$ has magnitude at most $\varepsilon$, and the unstable coordinate has magnitude at most $\varepsilon\lambda^{-(n-1)}$ by the bound at time $n-1$. The atom consequently has area at most
\[
 C_\varepsilon\lambda^{-n}.                             \tag{3}
\]
For atom probabilities $p_i$ satisfying (3),
\[
 H\!\left(\bigvee_{j=0}^{n-1}f_A^{-j}\mathcal P\right)
 =\sum_i p_i\log(1/p_i)
 \geq n\log\lambda-\log C_\varepsilon.
\]
Thus $h(f_A)\geq\log\lambda>0$, and (1) gives finiteness. An exact entropy evaluation is unnecessary for this example. It proves that smoothness and full smooth volume are compatible with positive finite entropy, while the preceding argument showed that one-dimensional models cannot do this.

\subsubsection{4. An exact symbolic model which is not a smooth solution}
Let $\Sigma=\{0,1\}^{\mathbb Z}$ with independent fair coordinates and the left shift $\sigma$. Off the null set of ambiguous binary expansions, define
\[
 x=\sum_{j\geq0}\frac{s_j}{2^{j+1}},\qquad
 y=\sum_{j\geq1}\frac{s_{-j}}{2^j}.
\]
The positive and negative coordinate halves are independent, so this map takes product probability to Lebesgue area on $[0,1)^2$ and has a measurable inverse modulo null sets. A direct digit calculation conjugates the shift to the baker transformation
\[
 B(x,y)=\left(2x-\lfloor2x\rfloor,
             \frac{y+\lfloor2x\rfloor}{2}\right).         \tag{4}
\]
Each branch has determinant one and maps a vertical half-rectangle onto a horizontal half-rectangle. The map is invertible almost everywhere: the branch is recovered from whether the second output coordinate lies below or above $1/2$.

The partition by the current symbol has $n$-fold join consisting of $2^n$ equiprobable atoms, so its entropy rate is $\log2$. It generates the two-sided symbolic sigma-algebra under all iterates, giving entropy $\log2$ for the shift and hence for (4). To see the upper bound directly, approximate a finite measurable partition in conditional entropy by a partition depending on finitely many coordinates. An $n$-join of a block depending on coordinates $-m,\ldots,m$ uses at most $n+2m$ independent bits, so its entropy divided by $n$ tends to at most $\log2$. The conditional-entropy error can be arbitrarily small. Thus the lower bound from the current-symbol partition is exact.

Nevertheless (4) is not a $C^\infty$ diffeomorphism of the square or of the torus obtained by identifying opposite sides. At an interior point of the cut $x=1/2$, the two limiting second coordinates are $y/2$ and $(y+1)/2$, differing by $1/2$ even modulo one. The first coordinates agree modulo one but the second do not.

Changing only the values on the cut cannot make the map continuous. More strongly, no continuous map can agree with this particular baker map almost everywhere on the torus. On each open branch, two continuous maps agreeing almost everywhere must agree everywhere: the set where they differ is open and would have positive area if nonempty. A putative continuous representative must therefore have the incompatible one-sided limits just computed. This is a genuine obstruction to that representative, not merely a derivative jump that could be ignored on a null set.

It does not show that the Bernoulli shift lacks some other smooth realization, and no such negative claim is intended. A different measurable conjugacy can change the geometry of all partition pieces. The lesson is that an exact piecewise affine symbolic model does not automatically become a smooth realization by ignoring its singular set.

\subsubsection{5. Matching entropy does not match the abstract transformation}
A second shortcut is to construct any smooth system having the same entropy as $T$. Even within the easiest realizable class this loses information. For irrational rotation $R_\alpha$, the eigenvalues of the Koopman operator on $L^2$ are precisely
\[
 \{e^{2\pi i k\alpha}:k\in\mathbb Z\}.
\]
Indeed the Fourier basis diagonalizes the operator, and an eigenfunction can have a nonzero coefficient only at frequencies with the asserted eigenvalue. A measurable isomorphism intertwines Koopman operators, so it preserves this eigenvalue set.

Take $\alpha=\sqrt2-1$ and $\beta=\sqrt3-1$. These rotations both have entropy zero, but $e^{2\pi i\alpha}$ is not an eigenvalue for $R_\beta$. Otherwise $\sqrt2-k\sqrt3$ would be an integer for some integer $k$. Squaring shows that if both integer coefficients are nonzero then $\sqrt3$ is rational; the cases with a zero coefficient are also impossible. Thus these two zero-entropy systems are not measurably isomorphic. Entropy is a necessary scalar invariant, not a complete prescription for the required model.

Embedding a symbolic system into a smooth map on an invariant Cantor set is also insufficient whenever the resulting invariant probability is singular to ambient volume. The target asks for the full smooth volume measure, so a model supported on a volume-zero invariant set does not meet it. Even when a construction produces a large invariant set, one must verify its measure and the actual isomorphism, rather than only its orbit coding.

\subsubsection{6. Verification and the first unresolved step}
The finite checks record $\det A=1$, the two eigenvalues, branchwise preservation and inversion of (4), and exact conjugacy for finite binary strings using a fixed zero-tail convention at dyadic endpoints. They test the explicit formulas only. No finite list of symbolic words can certify a smooth realization of an arbitrary transformation.

The proved partial results establish necessity of finite entropy, rigidity in dimension one, an ergodic positive-entropy smooth example, and two precise failures of shortcut constructions. What is still missing is a construction for a general finite-entropy ergodic $T$ that simultaneously gives a compact finite-dimensional manifold, convergence in every derivative to an invertible smooth map, preservation of a positive smooth volume form, and a measurable conjugacy of the entire system. None of the calculations supplies that construction. The universal realization problem remains unresolved by this attempt.

\subsection{Sources}
\begin{itemize}
\item[S1]M. Foreman, The complexity of the Structure and Classification of Dynamical Systems, v4 (2023).
\url{https://arxiv.org/html/2203.10655}.
\textit{Formulation source.}
\end{itemize}
\end{document}
